\section{Galileo 时空}\label{sec:galileo-spacetime}

所有物理事件构成的集合可以用连通的四维流形 $M$ 来模拟. 我们称 $M$ 为\textbf{时空}, 称其点为\textbf{事件}. 假定 $M$ 与 $\mathbb{R}^4$ 微分同胚, 且存在一类特殊的微分同胚 $x:M\to\mathbb{R}^4$, 称为\textbf{惯性参考系}. 一个惯性参考系给出全局坐标 $(x^0,x^1,x^2,x^3)=(t,x,y,z)$. 坐标 $t:M\to\mathbb{R}$ 称为与给定惯性参考系相联系的\textbf{时间函数}. 若 $t(p)=t(q)$, 则称两个事件 $p,q\in M$ 在该参考系中\textbf{同时}. 因此, 时间函数的水平集称为\textbf{同时超曲面}. 两个同时事件 $p,q\in M$ 之间的\textbf{距离}为
\[
  d(p,q)=\sqrt{\sum_{i=1}^3\bigl(x^i(p)-x^i(q)\bigr)^2}.
\]

粒子的运动用满足 $\mathrm{d}t(\dot c)\ne0$ 的光滑曲线 $c:I\to M$ 模拟. 一类特殊的运动是\textbf{自由粒子}的运动, 即不受任何外力的粒子的运动. 惯性参考系必须满足的特殊性质是, 自由粒子的运动总由直线表示. 换言之, 自由粒子相对于惯性参考系以恒定速度运动 (\textbf{Newton 惯性定律}). 特别地, 在惯性参考系中静止的粒子所作的运动, 是自由粒子的运动.

惯性参考系不是唯一的: 若 $x:M\to\mathbb{R}^4$ 是惯性参考系, $T:\mathbb{R}^4\to\mathbb{R}^4$ 是可逆仿射变换, 则 $T\circ x$ 是另一个惯性参考系. 事实上, 任意两个惯性参考系必由这种仿射变换联系起来 [参见 Exercise~\ref{ex:galileo-spacetime}(\ref{ex:straight-line-preserving-bijection-affine})].

作为 Newton 力学基础的 \textbf{Galileo 时空}, 是通过进一步要求各惯性参考系满足下列条件而得到的:
\begin{enumerate}[label=(\arabic*),leftmargin=*]
\item 对任意两个事件之间的时间间隔给出相同的结果 (从而对给定的两个事件是否同时也给出相同的判断);
\item 对同时事件之间的距离给出相同的结果.
\end{enumerate}

因此, 除去平移和反射, 惯性参考系之间的所有坐标变换都属于 \textbf{Galileo 群} $\operatorname{Gal}(4)$, 即保持时间函数及同时超曲面的 Euclidean 结构的保定向线性映射所成的群.

在分析只有一个空间维度重要的问题时, 可以使用更简单的二维 Galileo 时空. 若 $(t,x)$ 是与一个惯性参考系相联系的时空坐标, 且 $T\in\operatorname{Gal}(2)$ 是到全局坐标为 $(t',x')$ 的新惯性参考系的 Galileo 换基变换, 则
\begin{align*}
  \frac{\partial}{\partial t'}&:=T\left(\frac{\partial}{\partial t}\right)
  =\frac{\partial}{\partial t}+v\frac{\partial}{\partial x},\\
  \frac{\partial}{\partial x'}&:=T\left(\frac{\partial}{\partial x}\right)
  =\frac{\partial}{\partial x}
\end{align*}
其中 $v\in\mathbb{R}$, 因为必须有 $t=t'$, 从而
\[
  \mathrm{d}t\left(\frac{\partial}{\partial t'}\right)
  =\mathrm{d}t'\left(\frac{\partial}{\partial t'}\right)=1,
\]
而且我们要求保定向映射 $T$ 是同时超曲面 $\{t=0\}\equiv\{t'=0\}$ 的等距映射. 因此换基矩阵为
\[
  S=\begin{pmatrix}1&0\\v&1\end{pmatrix},
\]
其逆为
\[
  S^{-1}=\begin{pmatrix}1&0\\-v&1\end{pmatrix}.
\]

因此相应的坐标变换为
\[
  \begin{cases}
    t'=t\\
    x'=x-vt
  \end{cases}
  \qquad(v\in\mathbb{R})
\]
(\textbf{Galileo 变换}), 从而新参考系以速度 $v$ 相对于旧参考系运动 (因为曲线 $x'=0$ 就是曲线 $x=vt$). 注意, 只需改变 $v$ 的符号, 就可从 $S$ 得到 $S^{-1}$. 这符合预期, 因为旧参考系必以速度 $-v$ 相对于新参考系运动. 将这一事实称为\textbf{相对性原理}.

\begin{RiemannExercise}\label{ex:galileo-spacetime}
\begin{enumerate}[label=(\arabic*),ref=\arabic*,leftmargin=*]
\item\label{ex:lucas-transatlantic-ships}
(\emph{Lucas 问题}) 到 19 世纪末, Le Havre 与纽约之间已有定期的跨大西洋航线. 每日格林尼治时间 (GMT) 正午, Le Havre 和纽约各有一艘远洋轮船出发. 航行恰好需要七天, 因此抵达也是在格林尼治时间正午. 所以, 从 Le Havre 前往纽约的一艘远洋轮船, 会在出发时遇见一艘刚从纽约抵达的远洋轮船, 并在到达时遇见另一艘刚从纽约出发的轮船. 除此之外, 它还会遇见多少艘船? 分别在什么时间遇见? 这条航线总共需要多少艘船? (\textbf{提示}: 将各船的运动表示为二维 Galileo 时空中的曲线.)
\item\label{ex:free-particle-has-constant-velocity-in-inertial-frame}
验证自由粒子相对于惯性参考系以恒定速度运动.
\item\label{ex:straight-line-preserving-bijection-affine}
设 $f:\mathbb{R}^n\to\mathbb{R}^n$ ($n\geq2$) 是将直线映为直线的双射. 证明 $f$ 必是仿射函数, 即对所有 $x\in\mathbb{R}^n$ 有
\[
  f(x)=Ax+b,
\]
其中 $A\in GL(n,\mathbb{R})$, $b\in\mathbb{R}^n$.
\item\label{ex:galileo-group-as-euclidean-isometry-group}
证明 Galileo 群 $\operatorname{Gal}(4)$ 是 $GL(4,\mathbb{R})$ 中由形如
\[
  \begin{pmatrix}1&0\\v&R\end{pmatrix}
\]
的矩阵组成的子集, 其中 $v\in\mathbb{R}^3$, $R\in SO(3)$. 推出 $\operatorname{Gal}(4)$ 与 Euclidean 三维空间 $\mathbb{R}^3$ 的保定向等距映射群同构.
\item\label{ex:two-dimensional-galileo-subgroup}
证明 $\operatorname{Gal}(2)$ 是 $\operatorname{Gal}(4)$ 的子群.
\end{enumerate}
\end{RiemannExercise}
